%%%%%
%
\section{Calcul de spectre en concentrations par la m\'ethode \`a Duroure}

\itemstyle{Op\'erateur~:} {\sf Duroure\_2D}
\index{Duroure\_2D}

%%%%%OBJET%%%%
\itemstyle{Objet~:}
Cet op�rateur calcule le spectre moyen en concentrations, par la m�thode �
Duroure, des donn�es d'une sonde PMS 2D-C ou 2D-P.

%%%%%ARGUMENTS%%%%
\bigskip
%%%%%ENTREES%%%%
\begin{tabarg}
\multicolumn{3}{l}{\titretab {Entr\'ees}}\\
\hline
$Time_{reg}$ & Vecteur, {\tt Float\_64} & Temps avec fr\'equence de 1 Hertz ($s$)\\
$Image_{2D}$ & Matrice X1024, {\tt Int\_32} & Image\\
$Time_{asyn}$ & Vecteur, {\tt Float\_64} & Temps de la variable d'entr\'ee Image ($s$)\\
$TAS_{2D}$ & Vecteur, {\tt Int\_16} & True Air Speed ($m/s$)\\
$Code\_Hexa\_sonde$ & Coefficient, {\tt String} & Code hexadecimal de la sonde\\
$Masque\_sonde$ & Coefficient, {\tt String} & Masque hexadecimal des donn\'ees images\\
$Bits\_temps$ & Coefficient, {\tt Int\_16} & Nombre de bits de codage du temps dans les s\'eparateurs\\
$Nbval_{as}$ & Coefficient, {\tt Int\_32} & Nombre d'\'el\'ements dans le vecteur de temps ascynchrone\\
$Temps_{debut}$ & Coefficient, {\tt Int\_32} & Temps du d\'ebut de l'analyse en secondes depuis minuit\\
$Temps_{fin}$ & Coefficient, {\tt Int\_32} & Temps de la fin de l'analyse en secondes depuis minuit\\
$Tronq_{debut}$ & Coefficient, {\tt Int\_16} & Num\'ero de la diode \`a consid\'erer (d\'ebut)\\
$Tronq_{fin}$ & Coefficient, {\tt Int\_16} & Num\'ero de la diode \`a consid\'erer (fin)\\
$Sonde$ & Coefficient, {\tt String} & Nom de la sonde, soit 2D-P ou 2D-C\\
\end{tabarg}

\medskip
%%%%%SORTIES%%%%
\begin{tabarg}
\multicolumn{3}{l}{\titretab {Sorties}}\\
\hline
$Spec\_tot_{2D}$ & Vecteur, {\tt Float\_32} & Spectre total 2D\\
$Spec\_moy_{2D}$ & Matrice X30, {\tt Float\_32} & Spectre moyen total 2D\\
$Spec_{2D}$ & Matrice X30, {\tt Float\_32} & Spectre 2D\\
$Spec\_dmoy_{2D}$ & Vecteur, {\tt Float\_32} & Diametre moyen 2D\\
$Spec\_rugmoy_{2D}$ & Vecteur, {\tt Float\_32} & Rugosite moyenne 2D\\
\end{tabarg}


%%%%%FORMULE%%%%
\itemstyle{Formule}

%%% Duroure_2D


%%%%%COMMENTAIRE%%%%
\itemstyle{Commentaire~:}
Les donn�es images de la sonde PMS-2D-C ou 2D-P sont analys�es par imagette.
La m�thode de d�tection des imagettes a �t� modifi�e afin d'utiliser la m�me
que celle de l'op�rateur Compte\_img\_2D.
La taille d'une particule non-tronqu�e est estim�e en comptabilisant le nombre
de pixels actifs de l'imagette (la surface). Alors,
le diam�tre de la particule correspond
au diam�tre d'une particule circulaire ayant la m�me surface.
La taille d'une particule doublement tronqu�e est calcul�e correspond au
diam�tre selon l'axe des X.
Les particules simplement tronqu�es ne sont pas trait�es car il n'y a pas de
m�thode qui fonctionne bien pour calculer les concentrations avec celles-ci.
Les imagettes vides sont consid�r�es comme des particules de la plus petite
taille dans le spectre.
Certaines particules sont rejet�es, soient celles qui sont en splash
(dimension X au moins 4 fois plus grande que celle en Y) et les simplement
tronqu�es.
Le calcul des concentrations s'effectue en calculant un volume
d'�chantillonnage qui est calcul� en utilisant les valeurs fournies par PMS
pour la sonde utilis�e.
En sortie, deux fichier ASCII ayant les noms
Spec\_moy\_Duroure\_2d.dat et Spec\_moyt\_Duroure\_2d.dat sont cr��s et contiennent
le spectre moyen pour la p�riode sp�cifi�e, en unit�s de cm-3*micron-1.

%%%%%AUTEUR%%%%
\itemstyle{Auteur~:}
Christian Pag\'e CNRM/GMEI/MNP

%%%%%REFERENCE%%%%
\itemstyle{R\'ef\'erences~:}
